% ============================================================
% TP-1: The Truncated Photon and the Lattice Regularization
%       of Shutter-Induced Photon Creation
% Conscious Point Physics — Phenomena Series (quantum optics)
% Version history: TP-1/documentation_suite/changelog-TP-1.md (created at SHIP)
% ============================================================

\documentclass[11pt,letterpaper]{article}
\usepackage[utf8]{inputenc}
\usepackage[margin=1in]{geometry}
\usepackage[T1]{fontenc}
\usepackage{amsmath,amssymb,amsfonts,amsthm}
\usepackage{graphicx}
\usepackage{xcolor}
\usepackage{booktabs}
\usepackage{array}
\usepackage{enumitem}
\usepackage{float}
\usepackage[font=small, labelfont=bf]{caption}
\usepackage{parskip}
\usepackage{mdframed}
\usepackage[authoryear,round]{natbib}
\bibliographystyle{plainnat}
\usepackage{hyperref}
\hypersetup{colorlinks=true,linkcolor=blue,citecolor=blue,urlcolor=blue}

\newtheorem{theorem}{Theorem}[section]
\newtheorem{proposition}[theorem]{Proposition}
\newtheorem{lemma}[theorem]{Lemma}
\newtheorem{corollary}[theorem]{Corollary}
\newtheorem{definition}[theorem]{Definition}
\newtheorem{remark}[theorem]{Remark}

\newcommand{\lP}{l_P}
\newcommand{\tP}{t_P}
\newcommand{\EP}{E_P}
\newcommand{\SSV}{\mathrm{SSV}}
\newcommand{\phig}{\varphi}
\newcommand{\Tr}{\operatorname{Tr}}
% --- Paper-specific macros ---
\newcommand{\wg}{\omega_\gamma}
\newcommand{\wcut}{\omega_{\mathrm{cut}}}
\newcommand{\wP}{\omega_P}
\newcommand{\Nexp}{\langle N\rangle}

\title{\textbf{TP-1: The Truncated Photon and the Lattice
Regularization of Shutter-Induced Photon Creation}\\[6pt]
{\large Conscious Point Physics --- Phenomena Series (Quantum Optics)}\\[4pt]
{\Large \textbf{Version 0.1 --- 20 June 2026}}}

\author{Thomas Lee Abshier, ND \and Claude Opus (Anthropic)}

\date{\normalsize Hyperphysics Institute, Kalispell, Montana\\
\url{https://hyperphysics.com}\\
\href{mailto:drthomas007@protonmail.com}{drthomas007@protonmail.com}\\[4pt]
\small OSF DOI: [pending registration]}

\begin{document}
\maketitle

\begin{abstract}
\citet{RukanGullaSkaar2026} show, within standard quantum optics, that truncating a
single photon with an optical shutter does not yield a photon-or-vacuum mixture but a
superposition and mixture of photon numbers up to infinity, a state that is
nonetheless \emph{locally} equivalent to a single photon on one side of the shutter and
vacuum on the other. The expected photon number $\Nexp$ diverges only in the limit of an
instantaneously closing shutter; for realistic speeds even a thousand created photons is
extremely unlikely. We ask whether this effect is compatible with the postulates of
Conscious Point Physics (CPP), and we find that it is --- structurally, and in three
independent places. (i) A photon in CPP is a perturbation of the Dipole-Pair (DP) Sea, a
real propagating disturbance in a physical medium, so a \emph{cuttable} photon-wave is
native rather than paradoxical. (ii) Many-photon Fock content is ordinary bosonic
multi-occupation of the $600$-cell eigenmodes (QM-5 second quantization). (iii) The
local-simple / global-complex structure is the standard partial trace over DP-Sea
degrees of freedom under the Nexus global-consistency constraint (QM-4 measurement). The
creation mechanism the authors identify --- violation of time-translation invariance ---
maps onto a time-dependent reconfiguration of the DP-Sea boundary by the shutter, i.e.\ a
dynamical-Casimir-type process, reproducing their Noether reasoning rather than
conflicting with it. The non-trivial, framework-specific result is regularization: the
truncation divergence is \emph{ultraviolet} in origin, and CPP carries a hard cutoff at
$k_{\max}=\pi/\lP$ (the same cutoff that renders the electron self-energy finite in QM-5).
Therefore $\Nexp$ is finite for every physical shutter, and is bounded for \emph{all}
shutters because none can switch faster than one Absolute Moment $\tP$. Modeling the
created-photon spectrum with the mild (logarithmic) ultraviolet form
$dN/d\omega=\kappa/\omega$ --- the divergence class selected by the authors' own
``even a thousand photons is unlikely'' statement --- gives a Planck-set ceiling
$\Nexp_{\max}\simeq\kappa\,\ln(\wP/\wg)\approx 64\,\kappa$ for an optical photon: order
tens, finite, with $\kappa$ an $O(1)$ coefficient fixed only by importing the
Rukan--Gulla--Skaar Bogoliubov kernel and cutting its frequency integral at $\pi/\lP$
(registered as OPEN-TP-1). The honest status is therefore: \emph{compatible at the
structural level, with a framework-specific finiteness result that is robust and a
numerical ceiling whose coefficient is open; zero new axioms; no theorem claimed.}
\end{abstract}

\noindent\textbf{Keywords:} truncated photon; single-photon states; photon-number
distribution; dynamical Casimir effect; Bogoliubov transformation; ultraviolet cutoff;
Planck scale; Conscious Point Physics; DP Sea; 600-cell; locality in quantum field theory;
decoherence; Nexus; zero-parameter framework.

\medskip
\noindent\textbf{Plain Language Summary:} A photon is supposed to be indivisible, yet you
can clearly block part of a light pulse with a fast shutter. A recent study shows that
when you do this, the leftover light is not simply ``one photon or none'' --- it is a
strange blend of zero, one, two, and in principle endlessly many photons, even though an
observer on either side of the shutter would see only a single photon on one side and
nothing on the other. We show that this fits naturally into the Conscious Point Physics
picture, where a photon is a ripple in a real underlying medium and can be clipped like
any ripple. The one place the two pictures part ways is the ``endlessly many'': in the
standard theory the photon count blows up to infinity if the shutter could close
infinitely fast, whereas Conscious Point Physics has a smallest possible tick of time,
so the shutter can never be that fast and the photon count stays finite --- at most a few
dozen for ordinary light.

\bigskip
\raggedright
\tableofcontents
\newpage

% ============================================================
% SECTION 1: Introduction
% ============================================================
\section{Introduction}

A photon is an elementary excitation of the electromagnetic field and cannot be cut into
two photons. Yet one can plainly remove part of an optical pulse with a shutter.
\citet{RukanGullaSkaar2026} (hereafter RGS) resolve the apparent tension within ordinary
quantum optics. Because exact single-photon states have infinite temporal tails, a state
that ``looks like a single photon to the left of the shutter and vacuum to the right''
cannot be a one-photon state; RGS show that the truncated state is instead a
superposition and mixture of photon numbers from zero up to infinity, while remaining
\emph{locally} indistinguishable from a single photon or from vacuum in the two disjoint
regions. The creation of these photons is, in their analysis, a consequence of the
violation of time-translational invariance introduced by the shutter; by Noether's
theorem one then expects an apparent failure of energy conservation. The expected photon
number diverges only as the shutter closes infinitely fast; for realistic shutter speeds
even a thousand created photons would be extremely improbable.

This paper asks a single question: \emph{is the truncated-photon effect compatible with
the postulates of Conscious Point Physics?} The answer developed below is yes ---
compatible structurally, and in three independent places --- and the one framework-specific
consequence is a regularization of the RGS divergence by the CPP ultraviolet cutoff. We
do not re-derive the RGS result; we accept it as a correct statement about the continuum
theory (noting it is not uncontested) and read it against the CPP substrate.

The logic has a useful self-checking feature. CPP claims (QM-5, \citealp{CPP-QM5}) that
standard quantum field theory is recovered as the continuum limit of $600$-cell lattice
dynamics. Since the RGS result \emph{is} standard QFT, CPP must reproduce it in that
limit, including the divergence. Bare compatibility is therefore nearly automatic and is
not, by itself, the contribution. The contribution is the place where CPP departs from
the continuum: the lattice carries a hard ultraviolet cutoff, and the truncation
divergence is ultraviolet, so CPP renders $\Nexp$ finite for every physical shutter.

\subsection{Open Problems Addressed}
This paper opens the quantum-optics domain of \texttt{series\_phenomena/} with its first
member, the truncated-photon phenomenon, and registers OPEN-TP-1 (fix the $O(1)$
coefficient $\kappa$ of the Planck-set ceiling by importing the RGS Bogoliubov kernel and
cutting its frequency integral at $\pi/\lP$). It draws load-bearing structure from QM-5
(second quantization and the $\pi/\lP$ cutoff, \citealp{CPP-QM5}) and QM-4
(measurement as partial trace over the DP Sea under the Nexus, \citealp{CPP-QM4}), and
it adopts the framework-conditional, zero-new-axiom, no-theorem posture established by
EU-1 (\citealp{CPP-EU1}).

% ============================================================
% SECTION 2: The RGS result
% ============================================================
\section{The Truncated Photon in the Continuum}\label{sec:rgs}

Let a single-photon wavepacket be
\begin{equation}
|1_\xi\rangle = \int d\omega\,\xi(\omega)\,\hat a^\dagger(\omega)\,|0\rangle,
\qquad \int d\omega\,|\xi(\omega)|^2 = 1 ,
\end{equation}
with temporal envelope the Fourier transform of $\xi(\omega)$. The shutter is not a
number-conserving projector on the one-photon sector; it is a time-dependent boundary
condition imposed on the field. RGS show that the resulting state has support on every
photon-number sector,
\begin{equation}
|\psi_{\mathrm{cut}}\rangle = \sum_{n=0}^{\infty} c_n\,|n\rangle_{\mathrm{(mode\ structure)}},
\qquad \Nexp=\sum_n n\,|c_n|^2 ,
\end{equation}
and that the two reduced states obtained by tracing over the field on either side of the
shutter are, respectively, a single photon and vacuum to the precision set by a narrow
transition region. The hallmark of the mechanism is that $\Nexp$ grows without bound only
as the switching becomes instantaneous: a sharper shutter injects higher-frequency
content, and an idealized step demands content at arbitrarily high frequency. This is the
structural signature of an \emph{ultraviolet} divergence, and it is the hinge of
everything below.

\paragraph{Why this is dynamical-Casimir physics.} A shutter that changes the field's
boundary condition in time is, formally, the same class of process as a moving mirror or
a cavity of time-varying length \citep{Moore1970,Dodonov2020}, which converts vacuum and
single-quantum input into multi-photon output by a Bogoliubov transformation mixing
$\hat a$ and $\hat a^\dagger$. The dynamical Casimir effect has been observed in
superconducting circuits \citep{Wilson2011}. The number of created quanta is
$\Nexp=\sum_k|\beta_k|^2$ with $\beta_k$ the Bogoliubov coefficients of the time-dependent
boundary; instantaneous switching makes $\beta_k$ fall off slowly in $k$ and the sum
diverge, while smooth switching over a time $\tau_s$ suppresses $\beta_k$ above
$\omega\sim1/\tau_s$ and the sum converges. RGS's ``infinite only if instantaneous'' is
exactly this statement.

% ============================================================
% SECTION 3: The CPP compatibility map
% ============================================================
\section{Compatibility with the CPP Postulates}\label{sec:compat}

We show that each surprising feature of the truncated photon lands on a standing CPP
mechanism rather than requiring a new one.

\subsection{A photon is a cuttable disturbance of a real medium}
In CPP a photon is not a fundamental object but a perturbation of the Dipole-Pair (DP)
Sea --- the lattice vacuum whose ground-state dipoles transmit all interactions. A
propagating photon is a propagating disturbance \emph{in a medium}, and truncating a
disturbance in a medium is unremarkable. Where the continuum theory finds it awkward that
an ``indivisible'' particle can have part of it removed, CPP finds the extended,
clippable photon-wave to be the natural description: the wave aspect of wave--particle
duality is carried by an actual substrate excitation, and the shutter acts on that
excitation directly.

\subsection{Many-photon content is ordinary bosonic occupation}
QM-5 derives second quantization from the $600$-cell eigenmode expansion: any DI-bit
amplitude configuration is expanded as $\psi_i=\sum_k a_k u_k(i)$ over the lattice
eigenvectors $u_k$, and promoting $a_k\to\hat a_k$ with bosonic commutators gives the
field operator $\hat\phi(r_i)=\sum_k(\hat a_k u_k(i)+\hat a_k^\dagger u_k^*(i))$. Photon
modes are bosonic (no occupation exclusion), so a state spanning many photon numbers is
ordinary multi-occupation of Sea modes. The ``$0$ to $\infty$ mixture'' is a superposition
over occupation numbers of the relevant DP-Sea modes --- native, not exotic.

\subsection{Local-simple / global-complex is the measurement structure of QM-4}
QM-4 solves the CPP measurement problem by identifying apparent collapse with
decoherence: a local subsystem couples to a subset of DP-Sea modes, its reduced density
matrix $\rho=\Tr_{\mathrm{Sea}}|\psi\rangle\langle\psi|$ is simple, and the Nexus enforces
global unitarity so that the full state retains all structure. The truncated photon's
``single photon to the left, vacuum to the right, $0\!\to\!\infty$ globally'' is precisely
the system/environment partial trace applied \emph{across a spatial boundary}. The feature
RGS find strangest is, in CPP, the expected output of the standard decoherence account.

\subsection{The creation mechanism: a driven DP-Sea boundary}
RGS attribute photon creation to broken time-translation invariance. CPP states the same
physics mechanistically: the shutter is a macroscopic, time-dependent reconfiguration of
the DP-Sea boundary, doing work on the Sea --- a dynamical-Casimir-type excitation from a
driven boundary (\S\ref{sec:rgs}). The Nexus enforces global consistency at each Absolute
Moment, but it does not demand that an open, externally driven subsystem conserve energy:
the photon subsystem is not isolated, and the shutter is the external agent supplying the
energy. CPP therefore reproduces the Noether/energy-bookkeeping reasoning of RGS rather
than contradicting it.

% ============================================================
% SECTION 4: Regularization
% ============================================================
\section{Lattice Regularization of the Divergence}\label{sec:reg}

The one place CPP departs from the continuum is the fate of the divergence.

\begin{proposition}[Lattice regularization of shutter-induced photon creation]
\label{prop:reg}
Assume (a) the truncated-photon expected number $\Nexp$ diverges only in the
instantaneous-shutter limit, i.e.\ the divergence has purely ultraviolet origin
\citep{RukanGullaSkaar2026}; and (b) CPP imposes a hard ultraviolet cutoff
$k_{\max}=\pi/\lP$ on all field modes, equivalently a maximal angular frequency
$\wP=1/\tP$ \citep{CPP-QM5}. Then $\Nexp$ is finite for every physical shutter, and is
bounded above over the set of \emph{all} physical shutters, because no shutter can switch
faster than one Absolute Moment $\tP$.
\end{proposition}

\begin{proof}[Argument]
Write the created-photon spectrum as $dN/d\omega$, supported on $[\wg,\wcut]$ where $\wg$
is the carrier frequency and $\wcut$ the effective spectral cutoff set by the switching
sharpness ($\wcut\sim1/\tau_s$). By (a), $\Nexp=\int_{\wg}^{\wcut}(dN/d\omega)\,d\omega$
diverges only as $\wcut\to\infty$. By (b), every physical mode has $\omega\le\wP$, so
$\wcut\le\wP$ and $\Nexp\le\int_{\wg}^{\wP}(dN/d\omega)\,d\omega<\infty$. The bound is
uniform over shutters because the sharpest physical switching, one Absolute Moment, gives
$\wcut=\wP$.
\end{proof}

The argument uses only the \emph{origin} of the divergence (ultraviolet) and the
\emph{existence} of the CPP cutoff. It does not depend on the detailed spectrum, so the
finiteness conclusion is robust; the numerical value of the ceiling is not.

\paragraph{The numerical ceiling.} Take the mild, logarithmic ultraviolet form
$dN/d\omega=\kappa/\omega$, so $\Nexp(\wcut)=\kappa\ln(\wcut/\wg)$. This is the divergence
class \emph{selected by RGS's own qualitative statement}: logarithmic growth means
reaching $N$ photons requires $\wcut\sim\wg e^{N/\kappa}$, an exponentially sharp shutter,
which is why ``even a thousand photons would be extremely unlikely.'' A stronger
(power-law) divergence would make moderate photon numbers easy and is excluded by that
statement. With the CPP cutoff $\wcut\le\wP$,
\begin{equation}
\Nexp_{\max}=\kappa\,\ln\!\frac{\wP}{\wg}.
\end{equation}
For an optical photon ($\lambda=500\,$nm, $\wg=3.767\times10^{15}\,\mathrm{rad/s}$) and
$\wP=1/\tP=1.855\times10^{43}\,\mathrm{rad/s}$,
\begin{equation}
\ln\!\frac{\wP}{\wg}=63.8,\qquad \Nexp_{\max}\approx 64\,\kappa,
\end{equation}
an order-tens ceiling. The coefficient $\kappa$ is the only model-dependent input; fixing
it requires the RGS Bogoliubov kernel $B(\omega,\omega')$ evaluated with its frequency
integral cut at $\pi/\lP$ (OPEN-TP-1).

\paragraph{Sharper than ``unlikely.''} RGS say large photon numbers are improbable for
realistic shutters. CPP says more: for $\kappa$ of order $0.1$--$1$, even $\sim\!10^2$
created photons would require $\tau_s<\tP$, which the lattice forbids by construction.
What is ``extremely unlikely'' in the continuum becomes ``impossible past the Planck
ceiling'' on the lattice. This is the same conversion QM-5 performs for the electron
self-energy, where the lattice turns an infinite renormalization into a finite one; here
it turns an infinite expected photon number into a finite, Planck-set one.

% ============================================================
% SECTION 5: Physical Interpretation
% ============================================================
\section{Physical Interpretation}\label{sec:physical}

In the CPP picture the entire episode is a driven-boundary problem in the DP Sea. Before
the shutter acts, the photon is a coherent single-quantum disturbance riding the Sea, of
finite spatial extent because it was emitted a finite number of Absolute Moments ago and
propagates at exactly $c=\lP/\tP$, one edge per tick. The shutter is a macroscopic,
time-dependent boundary inserted into the Sea. Slamming it does work on the Sea and
stirs its zero-point dipoles, exciting additional Sea modes --- the ``created photons.''
How many are excited is set by how abruptly the boundary moves relative to the Sea's
fastest mode, the Zitterbewegung oscillation near $1/\tP$. An observer confined to one
side couples only to the Sea modes there and reads a simple reduced state; the full
$0\!\to\!\infty$ structure lives in the global Sea configuration that the Nexus keeps
consistent.

The RGS programme's stated goal --- compactly supported, finite-tail photons so that
interactions have a clean causal order rather than infinite tails implying interaction
since the infinite past --- is, in CPP, already the native situation. The fundamental
degrees of freedom are local lattice excitations with strictly finite propagation speed,
so a physical photon is always a finite-extent disturbance built over finitely many ticks.
The infinite tails are an artifact of the continuum Fourier idealization; the
compactly-supported photon the authors construct by hand is the default object on the
$600$-cell.

\subsection{CP/GP Signature at This Scale}

\textbf{Load-bearing axioms.} The compatibility map rests on A1 (Conscious Points and the
DP-Sea vacuum they compose), A2 (the $600$-cell topology supplying the eigenmode basis and,
through $z$ and $\lP$, the cutoff), A3 (DI-bit propagation at $c=\lP/\tP$, which makes
photons finite-extent and supplies the maximal frequency $\wP=1/\tP$), and A4 (the Nexus,
which enforces global consistency while permitting an open driven subsystem to exchange
energy). The regularization (Prop.~\ref{prop:reg}) is carried specifically by A2+A3
through the single fact $\wcut\le\wP=1/\tP$.

\textbf{Visible vs.\ smoothed discreteness.} The lattice discreteness is visible here in
exactly one quantity --- the cutoff $\wP=1/\tP$ that caps the photon-number ceiling --- and
is otherwise smoothed away: the bosonic mode algebra, the Bogoliubov kernel, and the
partial-trace structure are all continuum-faithful, because that is what QM-5's
continuum-limit claim guarantees. The single surviving lattice fingerprint is the
finiteness of $\Nexp$ and its $\ln(\wP/\wg)$ ceiling. This is the orders-of-magnitude
dilution made explicit: a Planck-scale discreteness leaves only a logarithm of the
Planck-to-optical ratio in the observable.

\textbf{Macroscopic shadow.} The conventional statement ``the truncated-photon number
diverges only for an instantaneous shutter'' is the macroscopic shadow of there being a
shortest tick $\tP$: in the continuum the instantaneous shutter is a legal idealization
and the divergence is real; on the lattice it is illegal, and the divergence is its
absence-of-a-cutoff shadow.

% ============================================================
% SECTION 6: Mapping
% ============================================================
\section{CPP-to-Conventional-Physics Mapping}\label{sec:mapping}

\begin{table}[H]
\centering
\caption{Structural correspondence between the continuum (RGS) account and the CPP
substrate account of the truncated photon.}
\begin{tabular}{p{0.30\textwidth} p{0.32\textwidth} p{0.30\textwidth}}
\toprule
\textbf{Feature} & \textbf{Continuum / RGS} & \textbf{CPP substrate} \\
\midrule
Photon & elementary field excitation, infinite tails & DP-Sea perturbation, finite extent (finite ticks since emission) \\
Shutter & time-dependent boundary breaking time-translation invariance & macroscopic driven boundary doing work on the DP Sea \\
Photon creation & Bogoliubov mixing of $\hat a,\hat a^\dagger$ & dynamical-Casimir excitation of Sea modes \\
$0\!\to\!\infty$ mixture & multi-photon Fock support & bosonic multi-occupation of $600$-cell modes (QM-5) \\
Local simple / global complex & partial trace gives single photon / vacuum & $\Tr_{\mathrm{Sea}}$ under the Nexus (QM-4) \\
Energy bookkeeping & apparent non-conservation (Noether) & open subsystem driven by the shutter; Nexus global-consistent \\
$\Nexp$ as shutter $\to$ instantaneous & $\to\infty$ (UV divergence) & finite; $\wcut\le\wP=1/\tP$ (Prop.~\ref{prop:reg}) \\
\bottomrule
\end{tabular}
\end{table}

The correspondence is structural, at the level of degrees of freedom and the symmetry
that is broken, not a literal identification: CPP operates below the perturbative-QFT
granularity, and the lattice content appears only in the cutoff.

% ============================================================
% SECTION 7: Falsifiers
% ============================================================
\section{Falsifiers}\label{sec:falsifiers}

\begin{enumerate}[leftmargin=2em]
\item \textbf{Coefficient falsifier (OPEN-TP-1).} If the RGS Bogoliubov kernel, cut at
$\pi/\lP$, yields a created-photon spectrum that is \emph{not} in the mild (at most
logarithmic) ultraviolet class, the order-tens ceiling fails and must be replaced by
whatever $\int_{\wg}^{\wP}(dN/d\omega)\,d\omega$ actually gives. The finiteness claim
(Prop.~\ref{prop:reg}) survives regardless; only the number $\approx 64\kappa$ is at risk.
\item \textbf{Compatibility falsifier.} If a careful CPP treatment of the truncated photon
required violating global unitarity at the level of the full Sea state (not merely the
open subsystem), the QM-4 measurement account would be contradicted and the compatibility
claim of \S\ref{sec:compat} would fail.
\item \textbf{Cutoff falsifier.} If the divergence were shown to be infrared rather than
ultraviolet --- i.e.\ to survive any high-frequency cutoff --- the lattice cutoff would not
regulate it and Prop.~\ref{prop:reg} would be vacuous. RGS's ``infinite only if
instantaneous'' is direct evidence against this, but a full spectral computation
(OPEN-TP-1) is the proper test.
\item \textbf{Empirical (in principle).} A measured truncated-photon number distribution
whose tail does not cut off below the bound implied by the fastest realizable shutter
would constrain $\kappa$ and, at extreme shutter speeds far beyond present capability,
could in principle distinguish a cutoff theory from a strictly continuum one.
\end{enumerate}

% ============================================================
% SECTION 8: Conclusion
% ============================================================
\section{Conclusion}\label{sec:conclusion}

The truncated-photon effect of Rukan, Gulla and Skaar is compatible with Conscious Point
Physics, and in three independent ways it is more than compatible: the cuttable
photon-wave, the many-photon Fock content, and the local-simple/global-complex structure
each instantiate a standing CPP mechanism (DP-Sea perturbation, bosonic mode occupation,
partial trace under the Nexus), and the authors' own creation mechanism --- broken
time-translation invariance --- is the CPP driven-boundary (dynamical-Casimir) process.
The single framework-specific result is that the divergence, being ultraviolet, is
regulated by the CPP cutoff $\wP=1/\tP$: the truncated-photon expected number is finite
for every physical shutter and bounded for all of them, with an optical ceiling of order
$64\kappa$ in the mild-divergence class the authors' phenomenology selects.

\subsection{Swarm-Validation Contribution}
This paper adds no new zero-parameter numerical correspondence to the cumulative swarm
tally; its contribution is one framework-conditional structural result (the lattice
regularization, Prop.~\ref{prop:reg}) and one in-principle-falsifiable order-of-magnitude
prediction (the $\sim\!64\kappa$ optical ceiling) carried by the unchanged $9$-axiom
stack, with the coefficient $\kappa$ deferred to OPEN-TP-1. Consistent with programme
convention, the running swarm total is unchanged; the paper's weight is that an
independently published, contested continuum result fits the CPP substrate without
adjustment and is sharpened by it.

\subsection{Problem Status After This Paper}\label{sec:problem_status}
\begin{itemize}[leftmargin=2em]
\item NEW: \texttt{series\_phenomena/quantum\_optics/} domain opened; truncated-photon
phenomenon registered as its first member at gated-conjecture maturity (this v0.1 draft).
\item NEW: OPEN-TP-1 registered --- fix $\kappa$ by importing the RGS Bogoliubov kernel and
cutting its frequency integral at $\pi/\lP$.
\item Prop.~\ref{prop:reg} (lattice regularization): stated and argued at framework-
conditional level; no global theorem claimed.
\item QM-4, QM-5: load-bearing, unchanged; this paper is an application, not a revision.
\end{itemize}

\section*{Acknowledgements}
The truncated-photon result is due to Rukan, Gulla and Skaar; this paper only reads it
against the CPP substrate and makes no claim on their derivation. Claude Opus (Anthropic)
contributed the compatibility analysis, the regularization argument, and the numerical
ceiling; Thomas Lee Abshier directed the work and the framework.

\bibliography{TP-1_refs}

\appendix
% ============================================================
% APPENDIX: Numerical Verification
% ============================================================
\section{Numerical Verification}\label{app:numerics}

The script
\texttt{series\_phenomena/quantum\_optics/photon\_truncation/scripts/1700\_truncation\_regularization.py}
evaluates the ceiling and the regularization curve. With $\tP=5.391247\times10^{-44}\,$s
and an optical carrier $\wg=3.767\times10^{15}\,$rad/s it returns
$\wP=1/\tP=1.855\times10^{43}\,$rad/s, $\ln(\wP/\wg)=63.76$, hence
$\Nexp_{\max}=63.8\,\kappa$. In the log model $\Nexp(\wcut)=\kappa\ln(\wcut/\wg)$ grows
from $\approx 10$ at $\wcut=10^{20}\,$rad/s to $63.8$ at the cutoff and diverges only as
$\wcut\to\infty$; reaching $N\gtrsim10^2$ created photons (for $\kappa\lesssim1$) requires
$\tau_s<\tP$, which the lattice forbids. A power-law spectrum ($p=0.5$) is included to
confirm that finiteness-under-cutoff is spectrum-independent while the magnitude is not.

\end{document}
